\section{\ours{}: reading the actions back}
\label{sec:improve}

\paragraph{Idea.}
Every metric in \cref{tab:metrics} except the judges needs the real video of the same actions, and none checks whether the motion shown could be performed by the robot.
\ours{} uses only what a simulator always has, the commanded actions: an inverse dynamics model reads the end-effector trajectory back from the generated video, and we measure how far it deviates from the commanded trajectory and whether the arm's inverse kinematics can execute it (\cref{fig:pipeline}).
A video that shows the commanded motion closes the cycle from actions to video and back with a small error.
We use NVIDIA Cosmos3-Nano~\cite{nvidia2026cosmos3} in its inverse dynamics mode for Bridge, the domain \texttt{bridge\_orig\_lerobot}, without any training of ours.

\paragraph{Decoding the model output.}
For each pair of consecutive frames the model returns 10 numbers.
Following the encoder of~\cite{nvidia2026framework} for this domain, they are the relative pose $T_i^{-1}T_{i+1}$ of the end-effector, given as the translation and the first two columns of the rotation matrix, followed by the commanded gripper state, each quantile-normalised as
\begin{equation}
  \hat a = 2\,\frac{a-q_{01}}{q_{99}-q_{01}} - 1,
  \label{eq:quantile}
\end{equation}
where $q_{01}$ and $q_{99}$ are the published 1st and 99th percentiles of each dimension.
We invert \cref{eq:quantile}, complete the rotation matrix with the cross product of its two columns, project it onto the nearest rotation, and chain the relative poses from the true first pose $T_0$, expressed in the flange frame of the encoder, which gives the trajectory $\hat p_{1:15}$ read from the video; nothing is fitted.
The Bridge states used by IRASim differ from the LeRobot version used by NVIDIA by a roll of $\pi$, with a median roll of 2.47\,rad in our clips against 0.00 in LeRobot, so we set $R_{\mathrm{LeRobot}} = R_{\mathrm{IRASim}}\, R_x(\pi)$.
Encoding our states with the code of~\cite{nvidia2026framework} and decoding them with ours returns them within $2\cdot10^{-4}$\,mm.

\paragraph{Cycle error.}
The cycle error of a clip is the mean distance between the end-effector positions read from the video and the commanded ones,
\begin{equation}
  E_{\mathrm{cycle}} = \frac{1}{15}\sum_{t=1}^{15} \lVert \hat p_t - p_t \rVert_2 ,
  \label{eq:idmscore}
\end{equation}
in centimetres, lower being better.

\paragraph{Executability.}
We solve the inverse kinematics of the WidowX 250 6-DOF arm, the Bridge robot, for each of the 16 poses read from the video, following the Interbotix driver: the screw axes and home pose of the arm, Newton iterations with tolerances of 1\,mrad and 1\,mm from three default initial guesses, the joint limits, and the joint velocity limit of $\pi$\,rad/s, checked between consecutive poses 0.2\,s apart at the 5\,Hz Bridge control rate.
Newton's method from the three default guesses misses valid solutions for low, steeply pitched poses, rejecting 4 of the 105 recorded Bridge trajectories, so we add a grid of 125 initial guesses with the same tolerances and limits.
A clip is executable when every pose has a solution within the joint limits and every step stays below the velocity limit.
With the wrong roll convention, none of the 384 recorded poses of the 24 core clips has a solution, so the check is sensitive to the frame convention.

\paragraph{Use for selection.}
Because \ours{} needs no real video, it can rank several samples generated from the same actions and keep the one with the lowest cycle error, which we test as best-of-$N$ selection.
